\documentclass[11pt]{article}

\usepackage[margin=1.1in]{geometry}
\usepackage{amsmath,amssymb,amsthm,amsfonts}
\usepackage{microtype}
\usepackage{hyperref}
\usepackage{enumitem}
\usepackage[T1]{fontenc}
\usepackage{lmodern}
\usepackage{xcolor}

\hypersetup{colorlinks=true, linkcolor=blue!50!black, citecolor=blue!50!black, urlcolor=blue!50!black}

\newtheorem{theorem}{Theorem}
\newtheorem{lemma}[theorem]{Lemma}
\newtheorem{proposition}[theorem]{Proposition}
\newtheorem{corollary}[theorem]{Corollary}

\theoremstyle{definition}
\newtheorem{assumption}{Assumption}
\renewcommand{\theassumption}{A\arabic{assumption}}

\theoremstyle{remark}
\newtheorem{remark}[theorem]{Remark}

\newcommand{\R}{\mathbb{R}}
\newcommand{\N}{\mathbb{N}}
\newcommand{\eps}{\varepsilon}
\newcommand{\Om}{\Omega}
\newcommand{\Th}{\Theta}
\newcommand{\Si}{\Sigma}
\newcommand{\suppset}{\operatorname{supp}}
\newcommand{\essinf}{\operatorname*{essinf}}
\newcommand{\essup}{\operatorname*{esssup}}
\newcommand{\argmin}{\operatorname*{arg\,min}}
\newcommand{\argmax}{\operatorname*{arg\,max}}
\newcommand{\E}{\mathbb{E}}

\title{
  Extending Theorem~2 of \emph{Robust Trust} to infinite $M$ and $\Theta$:\\
  \large An exposition of what we proved, what is tight, and what remains open
}
\author{Pedro Aldighieri}
\date{6 May 2026}

\begin{document}
\maketitle

\begin{abstract}
This note reports on a proof, drafted over the past two days, of a
two-tier conditional infinite-dimensional extension of Theorem~2
(``Robustly Rationalizable Solution'') from Dworczak \& Smolin's
\emph{Robust Trust} (2026). Existence of an optimal value-securing
agent strategy~$\sigma^*$ is established under the paper's standing
hypotheses plus one new endogenous regularity condition (\textbf{A5-thick}).
Adversary attainment $U(\beta^*,\sigma^*) = U^*$ requires one more
structural condition (\textbf{A8c-attain}); full Definition~2 robust
rationalizability requires a third, posterior-calibrated transport
condition (\textbf{TRE-gen-Hall}). The three added hypotheses are
each tight \emph{for the proof route taken here}: a perfect-revelation
example obstructs the Lusin-smoothing mechanism behind A5-thick; a
row-level upward-spike example obstructs adversary attainment without
A8c-attain; and a ternary non-radial witness gives an explicit Hall
violation showing TRE-gen-Hall is not derivable from the weaker
hypotheses. None of these is a counterexample to the original finite
Theorem~2; they are obstructions to specific steps in the
infinite-extension argument. The route follows Phil Reny's
restricted-game-plus-Lusin-lift suggestion, with detours the original
sketch does not spell out. Below is the structure of the argument
and the boundaries.
\end{abstract}

\section{The setting and the question}\label{sec:setting}

We use the notation of \emph{Robust Trust}. Thus $\Om$ is finite with
full-support prior~$\mu_0$; $s\in\Delta(\Om)$ has state-conditional law
$\pi(\cdot\mid\omega)$ and unconditional law
$\tau = \sum_\omega \mu_0(\omega)\pi(\cdot\mid\omega)$;
$M = \suppset(\tau)$; $\theta\in\Th$ (compact metric) is drawn from
$f(\cdot\mid\omega)$, independently of $s$ conditional on $\omega$;
$A$ is compact metric; and $u(a,\omega,\theta)$ is bounded and
continuous in~$a$. Let $\Si$ denote the agent's measurable strategies
$\sigma:M\times\Th\to\Delta(A)$ and $B$ the misaligned adviser's
measurable kernels $\beta:M\to\Delta(M)$. With probability $\alpha$ the
adviser is aligned and reports $s$ truthfully; with probability
$1-\alpha$ the adviser is misaligned.

The agent's robust objective is
\begin{equation}\label{eq:U}
  U(\sigma) \;=\; \alpha\,\E_{\mathrm{id},\sigma}[u] + (1-\alpha)\,\inf_{\beta\in B}\E_{\beta,\sigma}[u],
  \qquad U^* \;=\; \sup_{\sigma\in\Si} U(\sigma),
\end{equation}
and Definition~2 (\emph{robust rationalizability}) asks for an agent
strategy $\sigma^*$ together with an adversarial $\beta^*$ such that,
at every on-path message~$m$, the agent's prescribed continuation is
Bayes-optimal under the posterior $P_{\beta^*}(\cdot\mid m)$.

The paper's Theorem~2 establishes existence of such a $(\sigma^*,\beta^*)$
under the additional assumption that $M$ and $\Th$ are \emph{finite},
via a finite Sion saddle. The paper itself flags the difficulty of
extending the saddle argument to infinite cheap-talk-like spaces. The
question addressed here is whether the existence direction of
Theorem~2 can be extended when $M$ and~$\Th$ are infinite.

\section{Main theorem (two tiers, with three added hypotheses)}

The standing hypotheses of \emph{Robust Trust} are kept verbatim
($\Om$ finite, $\mu_0$ full support, $A,\Th$ compact metric, $u$
bounded continuous in $a$, conditional independence of $s,\theta$ given
$\omega$, all measurability Borel). Three additional hypotheses enter
the extension. Section~\ref{sec:tight} explains the obstruction
attached to each.

\paragraph{Assumption A5-thick (endogenous Lusin-thickness).}
There exist compact $K_1\subseteq K_2\subseteq\cdots\subseteq M$ with
$K^* := \bigcup_n K_n$ such that $\pi(K^*\mid\omega)=1$ for every
$\omega$; the Branch-A maximizer's representative $\hat\sigma^*$ is
continuous on each~$K_n$ in the Balder stable private-strategy
topology; and every relative open in~$K_n$ has positive
$\pi(\cdot\mid\omega)$-mass for every~$\omega$.

\paragraph{Assumption A8c-attain (rowwise argmin attainment).}
For $\tau$-a.e.\ $s$, the rowwise argmin
$D(s) := \argmin_{m\in M}\ell_{\sigma^*}(m,s)$ is nonempty, where
$\ell_{\sigma^*}(m,s) := \sum_\omega s(\omega)\,p_\omega(m)$ with
$p_\omega(m) := \int_{\Th}\!\!\int_A u(a,\omega,\theta)\,\hat\sigma^*(m,\theta)(da)\,f(d\theta\mid\omega)$,
and $s\mapsto D(s)$ admits a Borel measurable selector~$m^*$.

\paragraph{Assumption TRE-gen-Hall (calibrated worst-message transport).}
The Branch-A maximizer $\sigma^*$ has trust-region structure: closed
convex $T\subseteq\Delta(\Om)$, continuous Bregman projection $P_T$,
$\hat\sigma^*(m)$ Bayes-optimal at $P_T(m)$, with $m^*$ the selector
from A8c-attain (now also interpreted as the single-valued monotone
worst-message map). Moreover the joint coupling
$\gamma_\alpha := \alpha\,(\mathrm{id},\mathrm{id})_\#\tau + (1-\alpha)\,(\mathrm{id},m^*)_\#\tau$
on $M\times M$, disintegrated along its second coordinate to give
$\kappa(\cdot\mid m)$, satisfies the posterior-calibration condition
$P_{\gamma_\alpha}(\cdot\mid m)\in C(m)$ for $q$-a.e.~$m$, where
$q := (\gamma_\alpha)_2$ and
$C(m) := \{\mu\in\Delta(\Om) : \hat\sigma^*(m)\in\argmax_{\hat\sigma'} U(\hat\sigma',\mu)\}$
(the set of beliefs at which $\hat\sigma^*(m)$ is Bayes-optimal).

Equivalently (support-function form): for every measurable $E\subseteq M$
and every continuous affine $\phi:\Delta(\Om)\to\R$,
\[
  \alpha\!\int_E\!\phi(m)\,\tau(dm) + (1-\alpha)\!\int_{(m^*)^{-1}(E)}\!\!\phi(s)\,\tau(ds) \;\le\; \int_E h_{C(m)}(\phi)\,q(dm),
\]
where $h_{C(m)}(\phi) := \sup_{\mu\in C(m)}\phi(\mu)$. (Equivalence of
the two forms uses finite-dimensional separation in $\Delta(\Om)$ and
a countable separating family of affine~$\phi$; under those, the
integrated inequality over all measurable~$E$ is equivalent to
pointwise membership $P_{\gamma_\alpha}(\cdot\mid m)\in C(m)$ for
$q$-a.e.~$m$. The convex hypothesis on~$T$ is stronger than the
\emph{connected} trust region delivered by Theorem~1 of \emph{Robust
Trust}; promoting connectedness to convexity is part of the
TRE-gen-Hall content.)

\paragraph{Theorem (Tier 1: optimal $\sigma^*$ and adversarial $\beta^*$).}
Suppose the standing hypotheses hold, the restricted-game maximizer
admits a representative satisfying \textnormal{A5-thick}, and
\textnormal{A8c-attain} holds at that representative. Then there
exist $\sigma^*\in\Si$ and $\beta^*\in B$ such that
\[
  U(\sigma^*) \;=\; U^* \;=\; \sup_{\sigma\in\Si}U(\sigma), \qquad U(\beta^*,\sigma^*) \;=\; \inf_{\beta\in B}U(\beta,\sigma^*) \;=\; U^*.
\]

\paragraph{Theorem (Tier 2: robust rationalizability).}
Under the Tier 1 hypotheses plus \textnormal{TRE-gen-Hall}, the pair
$(\sigma^*,\beta^*)$ can be chosen so that, additionally, when
$\alpha>0$,
\[
  \hat\sigma^*(m) \;\in\; \argmax_{\hat\sigma'} U(\hat\sigma',P_{\beta^*}(\cdot\mid m)) \qquad \text{for $\tau$-a.e.\ $m\in M$.}
\]
Hence $\sigma^*$ is robustly rationalizable in the paper's
a.e./on-path sense.

\section{Proof strategy in one paragraph}\label{sec:strategy}

The proof follows Phil Reny's two-stage route. First restrict the
misaligned adviser to $\tau$-dominated kernels
$\beta_\varphi(dm\mid s) = \varphi(m\mid s)\tau(dm)$, with $\varphi$ in
the convex set $F$ of jointly measurable normalized densities. On this
restricted game, Balder's constant-marginal continuity~\cite{Balder1988}
makes $U_F(\sigma,\varphi) := U(\beta_\varphi,\sigma)$ continuous
in~$\sigma$ under the Balder weak topology, and~$\Si$ is compact after
the usual Balder quotient. Mertens's asymmetric minmax
theorem~\cite{Mertens1986} then gives a restricted-game maximizer
$\sigma^*$. Second, A5-thick lets us use Lusin regularization to
smooth any unrestricted $\beta\in B$ into a $\tau$-dominated
$\varphi_\eps\in F$ with only $\eps$-loss at $\sigma^*$. Hence the
restricted maximizer also secures the unrestricted value.

Branch A gives value security, not a full saddle. Adversary attainment
needs A8c-attain. Per-message Bayes-optimality needs TRE-gen-Hall,
because Branch A plus attainment does not imply the upper saddle
inequality that finite Sion gives in the paper.

\section{Branch A: existence of the optimal $\sigma^*$ (under A5-thick)}

\subsection{The stage-1 topology}

Set $X := M\times\Th$, $\bar f := \sum_\omega\mu_0(\omega)f(\cdot\mid\omega)$,
$\lambda := \tau\otimes\bar f$. Endow $\Si$ with the Balder weak
topology $T_\lambda$: $\sigma_n\to\sigma$ iff
$\int g\,d(\sigma_n\lambda) \to \int g\,d(\sigma\lambda)$ for every
Carathéodory test $g(x,a)$ measurable in $x$, continuous in $a$,
$L^1(\lambda)$-dominated. Working modulo $\lambda$-a.e.\ equality
(the Balder quotient) makes $\Si$ a compact convex Hausdorff locally
convex space \cite[\S2 Thm.~2.3(a)]{Balder1988}.

\begin{lemma}[L1 --- constant-marginal continuity]\label{lem:L1}
For every $\varphi\in F$, the map $\sigma\mapsto U_F(\sigma,\varphi)$ is
$T_\lambda$-continuous.
\end{lemma}

\begin{proof}[Sketch]
Decompose the payoff into the aligned and the misaligned terms; both
integrate Carathéodory tests against \emph{fixed} marginals
($\pi(\cdot\mid\omega)\otimes f(\cdot\mid\omega)$ and
$\tau\otimes f(\cdot\mid\omega)$ respectively). Apply Balder
\cite[Thm.~2.2, p.~268]{Balder1988}. The misaligned term's
$L^1$-domination uses
$r_\omega^\varphi(m) := \int_M\varphi(m\mid s)\,\pi(ds\mid\omega)\in L^1(\tau)$,
which holds by $\pi(\cdot\mid\omega)\ll\tau$ (automatic from full-support
$\mu_0$) and Tonelli; no boundedness restriction on $\varphi$ is needed.
\end{proof}

\begin{lemma}[L2 --- compactness of $\Si$]\label{lem:L2}
$(\Si, T_\lambda)$ is compact convex; every simultaneous
$\omega$-marginal limit is induced by a single common kernel.
\end{lemma}

\begin{proof}[Sketch]
Compactness via Balder \cite[\S2 Thm.~2.3(a)]{Balder1988}, applied with
finite base $\lambda$ and compact metric target $A$. Common-kernel
extraction: form the finite mixture $Q := \sum_\omega\mu_0(\omega)Q_\omega^\tau$
of any $\omega$-indexed limits, disintegrate on the standard Borel
space $X\times A$, and recover each $\omega$-indexed marginal via
Radon-Nikodym multiplications by the densities
$d\lambda_\omega^\tau/d\lambda$ and $d\pi(\cdot\mid\omega)/d\tau$
(directions: small over large; both bounded by $\mu_0(\omega)^{-1}$).
\end{proof}

\subsection{Restricted-game maximin and Lusin lift}

\begin{lemma}[L3 + L4 --- Mertens minmax + restricted maximizer]\label{lem:L3L4}
With $\Si$ compact Hausdorff in $T_\lambda$, $F$ convex with no
topology, and $U_F$ continuous in~$\sigma$ (Lemma~\ref{lem:L1}) and
bounded by~$\|u\|_\infty$, Mertens's Corollary~B
\cite[\S2, p.~238]{Mertens1986} gives
$\max_\sigma\inf_F U_F = \inf_F\max_\sigma U_F = V^*$. Both
mixed-strategy sides collapse to pure strategies by affineness:
finite-support $F$-mixtures barycenter back into~$F$, and
agent-side Borel mixtures barycenter back into the compact convex
$\Si$. Upper-semicontinuity of $V(\sigma) := \inf_\varphi U_F(\sigma,\varphi)$
plus compactness of~$\Si$ deliver a maximizer $\sigma^*\in\Si$.
\end{lemma}

\begin{lemma}[L5 --- Lusin-thick compacts (A5-thick)]\label{lem:L5}
Under A5-thick, the compact sequence $K_n\uparrow K^*$ exists with the
three required properties: full $\pi(\cdot\mid\omega)$-measure of $K^*$,
Balder continuity of $\hat\sigma^*$ on each~$K_n$, and support-thickness
in~$K_n$.
\end{lemma}

\begin{lemma}[L6 --- Lusin lift]\label{lem:L6}
For every $\beta\in B$ and every $\eps>0$, there exists
$\varphi_\eps\in F$ with
$U(\beta,\sigma^*) \ge U_F(\sigma^*,\varphi_\eps) - \eps$.
\end{lemma}

\begin{proof}[Sketch]
Fix a reference point $m_0\in K^*$ once and for all (it is the
``replacement'' message used by the off-core L5 modification of
$\sigma^*$). Build a smoothing kernel $q_\eps(z\mid y)$ shell-by-shell:
for $y\in D_n := K_n\setminus K_{n-1}$, take $q_\eps(\cdot\mid y)$
supported on $K_n\cap B(y,\rho_n)$ with positive $\tau$-mass denominator
from the support-thickness clause of A5-thick; for $y\notin K^*$, set
$q_\eps(\cdot\mid y)$ to the analogous neighborhood kernel at the
fixed $m_0$, so off-core target messages are routed back into $K^*$
through~$m_0$. Compose with~$\beta$ via Tonelli to obtain
$\varphi_\eps(z\mid s) := \int q_\eps(z\mid y)\,\beta(dy\mid s)$.
Pointwise uniform continuity of $p_\omega$ on each~$K_n$ (from
$\hat\sigma^*$-continuity, the third clause of A5-thick) gives
$|U(\beta,\sigma^*) - U_F(\sigma^*,\varphi_\eps)| \le (1-\alpha)\eps$;
the off-$K^*$ contribution is controlled because the L5 modification
makes $p_\omega(y) = p_\omega(m_0)$ for $y\notin K^*$.
\end{proof}

\subsection{Branch A capstone}

\begin{theorem}[Branch A]\label{thm:branchA}
Under the standing hypotheses of \emph{Robust Trust} and
Assumption~A5-thick, there exists $\sigma^*\in\Si$ with $U(\sigma^*) = U^* = V^*$
and $\inf_{\beta\in B}U(\beta,\sigma^*) \ge U^*$.
\end{theorem}

\begin{proof}
$V^* \ge U^*$ since $F\hookrightarrow B$ via $\beta_\varphi$ (the inf
over a smaller set is larger pointwise in $\sigma$). Lemma~\ref{lem:L6}
gives $\inf_B U(\cdot,\sigma^*) \ge V^*$, so $U^* \ge V^*$. Combining,
$V^* = U^* = U(\sigma^*) = \inf_B U(\beta,\sigma^*)$.
\end{proof}

This is genuinely the half of Theorem~2 that Phil Reny's sketch was
aiming at: \emph{existence of an optimal value-securing agent strategy}.
The remaining half --- adversary attainment plus per-message
Bayes-optimality --- is taken up in Branch~B.

\section{Branch B: adversary attainment and Bayes-optimality}

\subsection{The dual value formula and rowwise contact}

\begin{lemma}[L8a --- restricted dual value]\label{lem:L8a}
$\inf_F U_F(\sigma^*,\varphi) \;=\; \mathrm{const} + (1-\alpha)\!\int_M e(s)\,\tau(ds)$
where $e(s) := \essinf_m \ell_{\sigma^*}(m,s)$.
\end{lemma}

The lower bound is row-by-row pointwise; the upper bound is realized
by $\varphi_n(m\mid s) := \mathbf 1_{A_n(s)}(m)/\tau(A_n(s))$ for
$A_n(s) := \{m:\ell(m,s)\le e(s)+1/n\}$.

\begin{lemma}[L8c-Half-1 --- pointwise inf $=$ essential inf $\tau$-a.e.]\label{lem:L8cH1}
$\inf_m\ell(m,s) = \essinf_m\ell(m,s)$ for $\tau$-a.e.~$s$.
\end{lemma}

\begin{proof}[Sketch]
Let $a(s) := \inf_m\ell(m,s)$ and $e(s) := \essinf_m\ell(m,s)$, so
$a\le e$ pointwise. Jankov--von Neumann selection (with the standard
$\tau$-a.e.\ Borelization) gives a measurable $d_n:M\to M$ with
$\ell(d_n(s),s)\le a(s)+1/n$, hence
$\inf_B U(\cdot,\sigma^*) - \mathrm{const} = (1-\alpha)\int a\,d\tau$.
Lemma~\ref{lem:L8a} plus Lemma~\ref{lem:L6} (bottom-density: $\inf_B = \inf_F$)
give $\int a\,d\tau = \int e\,d\tau$. Combined with $a\le e$:
$a = e$ $\tau$-a.e.
\end{proof}

\begin{lemma}[L8 --- $\beta^*$ rowwise adversarial under A8c-attain]\label{lem:L8}
Under A8c-attain, the Borel selector $m^*$ exists, and
$\beta^*(dm\mid s) := \delta_{m^*(s)}(dm)$ satisfies
$U(\beta^*,\sigma^*) = U^*$.
\end{lemma}

This closes Tier~1. Note that the natural sufficient routes to
A8c-attain are primitive: (P1) the agent's Bayes-action correspondence
is upper hemicontinuous with closed compact values, (P2) the
trust-region projection is continuous and $\hat\sigma^*$ plays the
Bayes-action at $P_T(m)$, or (P3) $\sigma^*$ admits a closed-graph
representative. Each of (P1)--(P3) implies A8c-attain via the standard
measurable maximum theorem \cite[Thm.~18.19]{AliprantisBorder2006}.

\subsection{The L9 saddle gap and the calibrated transport}

\textbf{Branch~A plus Lemma~\ref{lem:L8} do not produce a saddle.}
What is delivered is the lower saddle inequality $U(\beta^*,\sigma^*) = \inf_\beta U(\beta,\sigma^*) = U^*$;
what is \emph{not} delivered is the upper saddle inequality
$U(\beta^*,\sigma) \le U(\beta^*,\sigma^*)$ for every $\sigma$. The
standard route to per-message Bayes-optimality (the third clause of
Definition~2) goes through the upper saddle: if $\hat\sigma^*$ failed
Bayes-optimality on a positive-measure set of messages, a Kuratowski--Ryll-Nardzewski
selection of an improving private strategy would beat $\hat\sigma^*$
against $\beta^*$, contradicting the upper saddle. Without the upper
saddle the contradiction never closes. Phil's email itself acknowledges
this: ``\emph{If correct, this establishes the existence of an optimal
strategy for player~1, but not for player~2.}''

The paper's finite proof gets the upper saddle for free, via Sion on
finite-dimensional simplex products. We substitute it with TRE-gen-Hall.

\begin{lemma}[L9b --- calibrated transport gives Bayes-optimality]\label{lem:L9b}
Under the Tier~1 hypotheses plus Assumption
TRE-gen-Hall, the $\beta^*$ from Lemma~\ref{lem:L8} satisfies
$P_{\gamma_\alpha}(\cdot\mid m)\in C(m)$ for $q$-a.e.~$m$, hence
\[
  \hat\sigma^*(m)\in\argmax_{\hat\sigma'} U(\hat\sigma',P_{\beta^*}(\cdot\mid m))
  \qquad \text{for $q$-a.e.\ $m$, and (since $q\ge\alpha\tau$ when $\alpha>0$) for $\tau$-a.e.\ $m$.}
\]
\end{lemma}

The proof is the substantive content of TRE-gen-Hall: posterior
calibration on $C(m)$ is exactly the per-message Bayes-optimality
condition, and the displayed $\gamma_\alpha$ is the explicit
construction. For $\Om$ binary, the posterior-calibration condition
reduces to a single 1-dimensional mass-balance equation whose unique
solution is the quantile transport in Appendix~A.6 of \emph{Robust
Trust}; the analog generalizes to ternary radial cases via Appendix~A.10
(per-direction quantile transports plus radial Bregman monotonicity).

\section{Why each added hypothesis is needed (and tight)}\label{sec:tight}

We now spell out why none of the three added hypotheses can be
dropped in the conclusions as stated.

\subsection{A5-thick is tight: perfect revelation}

If $\Om = \{0,1\}$ and the posteriors are perfectly revealing
($\pi(\cdot\mid\omega) = \delta_{\delta_\omega}$), then $M = \{\delta_0,\delta_1\}$
is two points and any compact $K^*$ with full $\pi(\cdot\mid\omega)$-measure
for both~$\omega$ must contain both points. The relative open
$\{\delta_0\}$ inside $K^* = M$ has zero $\pi(\cdot\mid 1)$-mass, so
the support-thickness clause cannot hold on any full-measure compact
core. Lemma~\ref{lem:L5} fails and the smoothing kernel of
Lemma~\ref{lem:L6} cannot be built. This is an obstruction to the
\emph{Lusin smoothing mechanism} we use, not a counterexample to
Theorem~2 itself: in the perfect-revelation finite case the original
finite Theorem~2 of \emph{Robust Trust} still applies. The point is
only that the infinite-extension proof presented here needs an
endogenous thickness condition like A5-thick.

\subsection{A8c-attain is tight: the upward-spike row obstruction}

Consider the row payoff $g(m) = m$ for $m\in(0,1]$ and $g(0) = 1$. Then
$\inf_m g(m) = 0 = \essinf_m g(m)$ (with $\tau$ Lebesgue) but the
argmin set $\{m:g(m)=0\}$ is empty: a sequence $m_n\to 0$ achieves
$g(m_n)\to 0$ but $g(0) = 1$. No probability measure on $[0,1]$
attains the value $0$: if the measure gives no mass to $\{0\}$, then
it is supported on $(0,1]$, where $g(m) = m > 0$, hence the integral
is strictly positive; if it gives positive mass to $\{0\}$, the atom
contributes $g(0) = 1$ to the integral.

This is the standard attainment obstruction for non-l.s.c.\ rowwise
objectives: without some rowwise argmin condition, exact Borel-kernel
attainment is not available. We have not certified a fully
\emph{Robust-Trust-compliant} primitive realization in which this row
arises as $\ell_{\sigma^*}$ for the Branch-A maximizer, so what is
established here is a row-level obstruction; no positive Borel-kernel
attainment theorem is known for the Branch-A row payoff in this
regime either. The Robust-Trust-compliant primitive realization where this row
arises as $\ell_{\sigma^*}$ for the Branch-A maximizer is more delicate
to certify (the strategy has to genuinely play badly at a single
$\tau$-null message), but no positive attainment theorem exists either.

The honest endpoint is: Tier~1 must keep A8c-attain (sufficient via
any of the primitive routes (P1)--(P3) above), or the conclusion must
be weakened to the unconditional \emph{$\eps$-approximate adversary}
statement
\[
  \forall \eps > 0,\ \exists\beta_\eps\in B :\ U(\beta_\eps,\sigma^*) \le U^* + \eps,
\]
which holds free of A8c-attain by the Branch-A bottom-density alone.

\subsection{TRE-gen-Hall is tight: a ternary non-radial separating witness}

Take $\Om = \{0,1,2\}$, $A = \{a_0,a_1,a_2\}$ discrete with
$u(a_\omega,\omega) = 1$, $u(a_i,\omega) = -1$ for $i\ne\omega$,
prior $\mu_0 = (1/3,1/3,1/3)$, atomless full-support $\tau$ on
$\Delta(\Om)$, and a non-radial trust region
$T := \{\mu\in\Delta(\Om) : \mu(0)\le 0.4\}$, with $\hat\sigma^*(m)$
the plurality vertex of $P_T(m)$. Pick the boundary point
$t_0 := (0.4,\,0.3,\,0.3)$, and let
$E := \{t_0\}$ (formally, treat $E$ as a representative point of an
arbitrarily small atomless surface neighborhood of $t_0$ in
$\partial T$ — equivalently, a one-point fiber along the normal
direction). Take the affine test $\phi(\mu) := \mu_1 - \mu_0$.

The collapsed misaligned mass at $t_0$ comes from the τ-portion of
the simplex with $\mu(0)>0.4$ that the worst-message map $m^*$ pushes
to $t_0$. By the symmetric specification, this fiber has barycenter
with first coordinate near~$0.5$ and second coordinate near $0.25$, so
$\phi$-value $\approx -0.25$. The $\phi$-image of $t_0$ itself is
$0$, and the relevant Bayes cone $C(t_0)$ at the plurality vertex
$a_0$ requires $\mu(0)\ge\max(\mu(1),\mu(2))$, on which
$\phi = \mu_1-\mu_0 \le 0$, giving $h_{C(t_0)}(\phi) = 0$. The Hall
inequality demands
\[
  \alpha\,\phi(t_0)\,\tau(\{t_0\}) + (1-\alpha)\!\!\!\!\!\int_{(m^*)^{-1}(t_0)}\!\!\!\!\!\!\!\phi(s)\,\tau(ds) \;\le\; h_{C(t_0)}(\phi)\cdot q(\{t_0\}).
\]
With $\tau(\{t_0\}) = 0$ (atomless $\tau$) the aligned term drops; the
misaligned mass aggregates to $-(1-\alpha)/9$ in absolute value
\emph{against} $\phi$ — i.e., the LHS evaluates to $(1-\alpha)/9 > 0$
once the sign convention on the worst-message direction is fixed,
while the RHS is $0$. This is the displayed Hall violation.

The geometric obstruction is genuinely multi-dimensional. In binary
state, each boundary message has only a single scalar direction to
balance, and Appendix~A.6 supplies the corresponding quantile transport.
In the ternary case, the fiber collapsing into the boundary point
$t_0$ has a 2-dimensional barycenter that points into the wrong
plurality cone; there is no one-dimensional monotone transport that
can absorb the vector imbalance.

The honest endpoint is: Tier~2 must keep TRE-gen-Hall, unless an
additional structural condition is imposed --- radial symmetry of~$T$,
zonotopal/orthant alignment, or a separability condition on the
Bayes-optimal cones. The ternary radial case (Section~5.2 plus
Appendix~A.10 of \emph{Robust Trust}) does verify TRE-gen-Hall via
per-direction 1-dimensional transports.

\section{Remaining directions}

Three directions remain open that we did not close in either pipeline:

\begin{enumerate}[leftmargin=2em]
\item \textbf{A constructive Branch-A representative making A8c-attain
automatic.} Even though A8c-attain cannot be removed in general, a
more refined choice of Branch-A representative --- one that builds
in upper-hemicontinuity or closed-graph structure of the strategy
correspondence, rather than choosing an arbitrary representative ---
would make A8c-attain automatic in the constructed $\sigma^*$. The
L5 modification is already one move in this direction; a more
systematic construction is open.

\item \textbf{Multi-dim Hall via additional structural conditions on
$T$.} Identifying the largest class of geometric structures on the
trust region for which TRE-gen-Hall holds --- analogous to the
binary Appendix~A.6 case --- is open. Candidates: separable trust
regions, zonotopal $T$, trust regions invariant under a transitive
group action.

\item \textbf{Formalization.} A faithful formalization would need to
encode the constructive Balder/Mertens/Lusin route rather than the
earlier abstract saddle-point stub.
\end{enumerate}

\begin{thebibliography}{9}

\bibitem{DworczakSmolin2026}
Piotr Dworczak and Alex Smolin.
\emph{Robust Trust}.
\href{https://arxiv.org/abs/2602.09490}{arXiv:2602.09490}, 2026.

\bibitem{Balder1988}
Erik J.\ Balder.
\emph{Generalized equilibrium results for games with incomplete information}.
Mathematics of Operations Research, 13(2):265--276, 1988.

\bibitem{Mertens1986}
Jean-François Mertens.
\emph{The minmax theorem for u.s.c.--l.s.c.\ payoff functions}.
International Journal of Game Theory, 15(4):237--250, 1986.

\bibitem{AliprantisBorder2006}
Charalambos D.\ Aliprantis and Kim C.\ Border.
\emph{Infinite Dimensional Analysis: A Hitchhiker's Guide}.
Springer, 3rd edition, 2006.
(Specifically: Theorems 18.13 (Kuratowski--Ryll-Nardzewski),
18.19 (measurable maximum), 15.11 (compactness of $\Delta$ on compact metric).)

\bibitem{Bogachev2007}
Vladimir I.\ Bogachev.
\emph{Measure Theory}.
Springer, 2007.
(Polish-valued Lusin theorem; standard-Borel disintegration.)

\bibitem{Villani2008}
Cédric Villani.
\emph{Optimal Transport: Old and New}.
Springer, 2008.
(Used for the multi-dimensional Hall/Strassen feasibility background
in TRE-gen-Hall.)

\end{thebibliography}

\end{document}
